% Source: https://artofproblemsolving.com/community/c3h1658901_geometry
% Stable identity: aops-topic-1658901-q01
The tight belt consists of three straight pieces, one alongside each side of the triangle, and three arcs, one around each circle.

\begin{asy}
usepackage("amsmath");
size(7cm);
defaultpen(fontsize(10pt));
pair A=(0,0),B=(12,0),C=(201/24,sqrt(121-(201/24)^2));real r=2;
pair[] V={A,B,C};
draw(A--B--C--cycle,dashed);
for(int i=0;i<3;++i){
  pair P=V[i],Q=V[(i+1)%3],R=V[(i+2)%3];
  pair n1=unit(rotate(-90)*(Q-P)),n2=unit(rotate(-90)*(R-Q));
  draw(P+r*n1--Q+r*n1,blue+linewidth(1));
  real a1=degrees(n1),a2=degrees(n2);if(a2<a1)a2+=360;
  draw(arc(Q,r,a1,a2),blue+linewidth(1));
  draw(circle(P,r));dot(P);
}
draw(B--B+r*dir(-30));label("$2$",B+0.5*r*dir(-30),NE,fontsize(9pt));
label("$A$",A,SW);label("$B$",B,SW);label("$C$",C,N);
label("$12$",(A+B)/2,N);label("$8$",(B+C)/2,SW);label("$11$",(A+C)/2,SE);
label("\mbox{}",(min(currentpicture,user=true)+max(currentpicture,user=true))/2,black); // keeps the text after the figure black when LaTeX typesets the labels inline
\end{asy}

\textbf{Straight pieces.} Each straight piece is a common outer tangent of two of the circles. The radii to its two points of tangency are perpendicular to it and both have length $2$, so the piece and the side of the triangle form a rectangle: the piece has the same length as the side. Together the straight pieces have length $8+11+12=31$.

\textbf{Arcs.} At a vertex with interior angle $\alpha$, the belt turns from the outward normal of one side to the outward normal of the next, through an angle of $360^\circ-90^\circ-90^\circ-\alpha=180^\circ-\alpha$. The three arcs turn through
\[
(180^\circ-\alpha)+(180^\circ-\beta)+(180^\circ-\gamma)=540^\circ-180^\circ=360^\circ
\]
in total, so together they make one full circle of radius $2$, of length $4\pi$.

The belt has length $\boxed{31+4\pi}$.
